\documentclass[a4paper,12pt]{article}

\usepackage[T2A]{fontenc}
\usepackage[utf8]{inputenc}
\usepackage[russian]{babel}

\usepackage{amsmath,amssymb,amsthm,mathtools}
\usepackage{helvet}
\renewcommand{\familydefault}{\sfdefault}
\usepackage{icomma}

\usepackage[a4paper,margin=2.2cm,headheight=16pt]{geometry}
\usepackage{microtype}

\usepackage{tikz}
\usetikzlibrary{arrows.meta,calc,positioning}

\usepackage{tkz-euclide}

\usepackage{pgfplots}
\pgfplotsset{compat=1.18}

\usepackage{tabularx,booktabs,longtable}
\usepackage{enumitem}
\usepackage{hyperref}
\usepackage{parskip}
\usepackage{fancyhdr}
\usepackage{setspace}
\usepackage{xcolor}

% -------------------------------------------------
% Палитра
% -------------------------------------------------
\definecolor{accent}{HTML}{008080}
\definecolor{darkaccent}{HTML}{004D4D}
\definecolor{textgray}{HTML}{333333}
\definecolor{softgray}{HTML}{EEF3F3}


% -------------------------------------------------
% Общие настройки
% -------------------------------------------------
\onehalfspacing
\setlength{\parindent}{0pt}
\setlength{\parskip}{0.55em}
\setlist[itemize]{leftmargin=1.7em,itemsep=0.25em,topsep=0.25em}
\setlist[enumerate]{leftmargin=1.9em,itemsep=0.45em,topsep=0.35em}

\hypersetup{
    colorlinks=true,
    linkcolor=darkaccent,
    urlcolor=darkaccent,
    pdftitle={Чек-лист: тригонометрические уравнения, теорема Безу и логарифмические неравенства},
    pdfauthor={Лёвин Артём Александрович}
}

\pagestyle{fancy}
\fancyhf{}
\fancyhead[L]{\small\color{textgray}ЕГЭ по математике}
\fancyhead[R]{\small\color{textgray}Николь Саркисьянц}
\fancyfoot[C]{\small\color{textgray}\thepage}
\renewcommand{\headrulewidth}{0.4pt}
\renewcommand{\headrule}{\hbox to\headwidth{\color{linegray}\leaders\hrule height \headrulewidth\hfill}}

\DeclareMathOperator{\tg}{tg}
\DeclareMathOperator{\ctg}{ctg}
\DeclareMathOperator{\arctg}{arctg}

\newcommand{\sectionline}{%
    \vspace{0.2em}
    \noindent\color{linegray}\rule{\linewidth}{0.6pt}
    \vspace{0.4em}
}

\newcommand{\checkitem}{\(\square\)\ }

\begin{document}

% =================================================
% ТИТУЛЬНЫЙ ЛИСТ
% =================================================
\begin{titlepage}
\thispagestyle{empty}

\begin{center}

\vspace*{2.2cm}

{\Large\color{darkaccent}\textbf{Чек-лист}}

\vspace{0.8cm}

{\huge\bfseries
Тригонометрические уравнения,\\[0.25em]
теорема Безу\\[0.25em]
и логарифмические неравенства
}

\vspace{1.1cm}

{\large
Комплексное повторение заданий\\
повышенного уровня ЕГЭ
}

\vfill

{\large\bfseries
Лёвин Артём Александрович\\
эксклюзивно для Николь Саркисьянц
}

\vspace{0.8cm}

{\large \today}

\vspace{2.2cm}

\end{center}

% Сдержанная кривая Безье внизу титульного листа
\begin{tikzpicture}[remember picture,overlay]
    \draw[
        color=accent,
        line width=1.2pt
    ]
    ([xshift=1.4cm,yshift=1.55cm]current page.south west)
    .. controls
    ([xshift=5.2cm,yshift=2.20cm]current page.south west)
    and
    ([xshift=-5.2cm,yshift=0.95cm]current page.south east)
    ..
    ([xshift=-1.4cm,yshift=1.55cm]current page.south east);
\end{tikzpicture}

\end{titlepage}

% =================================================
% ОСНОВНОЙ ЧЕК-ЛИСТ
% =================================================
\section*{\color{darkaccent}1. Главная стратегия занятия}
\sectionline

На занятии повторялись три взаимосвязанных навыка повышенного уровня:

\begin{itemize}
    \item преобразование тригонометрического уравнения к алгебраическому;
    \item разложение кубического многочлена с помощью корня и теоремы Безу;
    \item решение логарифмического неравенства с обязательным учётом ОДЗ;
    \item отбор корней тригонометрического уравнения на заданном промежутке.
\end{itemize}

\textbf{Главный принцип:} сначала привести задачу к стандартному виду, затем применять соответствующий алгоритм.

% -------------------------------------------------
\section*{\color{darkaccent}2. Тригонометрическое уравнение: унификация функций}
\sectionline

\checkitem Убедитесь, что во всех тригонометрических функциях стоит один и тот же аргумент.

\checkitem Посмотрите, можно ли выразить одну функцию через другую.

Основное тождество:
\[
\sin^2 x+\cos^2 x=1.
\]

Отсюда:
\[
\sin^2 x=1-\cos^2 x,
\qquad
\cos^2 x=1-\sin^2 x.
\]

Если в уравнении преобладает \(\cos x\), удобно выразить через него \(\sin^2x\).

\subsection*{\color{accent}Ключевой пример}

Решить:
\[
2\cos^2x+\sin^2x=2\cos^3x.
\]

Используем
\[
\sin^2x=1-\cos^2x.
\]

Тогда
\[
2\cos^2x+1-\cos^2x=2\cos^3x,
\]
то есть
\[
2\cos^3x-\cos^2x-1=0.
\]

Вводим замену
\[
t=\cos x.
\]

Обязательно фиксируем ограничение:
\[
-1\le t\le1.
\]

Получаем:
\[
2t^3-t^2-1=0.
\]

% -------------------------------------------------
\section*{\color{darkaccent}3. Замена \(t=\cos x\): что нельзя забывать}
\sectionline

При замене
\[
t=\cos x
\]
всегда выполняется
\[
t\in[-1;1].
\]

При замене
\[
t=\sin x
\]
также:
\[
t\in[-1;1].
\]

Это ограничение необходимо учитывать после решения алгебраического уравнения.

Например, если после преобразований получилось
\[
t=2,
\]
такой корень отбрасывается, поскольку
\[
2\notin[-1;1].
\]

% -------------------------------------------------
\section*{\color{darkaccent}4. Теорема Безу и поиск корня многочлена}
\sectionline

Для многочлена \(P(t)\):

\[
P(a)=0
\quad\Longleftrightarrow\quad
P(t)\ \text{делится на}\ t-a.
\]

Иными словами, если число \(a\) является корнем многочлена, то
\[
t-a
\]
является его множителем.

Для многочлена с целыми коэффициентами любой \textbf{целый} корень является делителем свободного члена.

Например, для
\[
2t^3-t^2-1
\]
целые кандидаты:
\[
t=\pm1.
\]

Проверяем:
\[
P(1)=2-1-1=0.
\]

Следовательно,
\[
t=1
\]
--- корень, а многочлен делится на
\[
t-1.
\]

\subsection*{\color{accent}Важное различие}

Теорема Безу отвечает на вопрос:

\[
P(a)=0
\quad\Rightarrow\quad
t-a\ \text{--- множитель}.
\]

Правило о делителях свободного члена помогает подобрать кандидатов на целый корень.

% -------------------------------------------------
\section*{\color{darkaccent}5. Деление многочлена уголком}
\sectionline

Разделим
\[
2t^3-t^2-1
\]
на
\[
t-1.
\]

Полезно явно записывать отсутствующий член:
\[
2t^3-t^2+0t-1.
\]

Последовательно получаем:

\[
\frac{2t^3-t^2+0t-1}{t-1}
=
2t^2+t+1.
\]

Следовательно,
\[
2t^3-t^2-1
=
(t-1)(2t^2+t+1).
\]

Уравнение принимает вид
\[
(t-1)(2t^2+t+1)=0.
\]

Первая скобка:
\[
t-1=0
\quad\Rightarrow\quad
t=1.
\]

Для второй:
\[
2t^2+t+1=0.
\]

Дискриминант:
\[
D=1-8=-7<0.
\]

Действительных корней нет.

Итак,
\[
t=1.
\]

% -------------------------------------------------
\section*{\color{darkaccent}6. Обратная замена}
\sectionline

Возвращаемся к
\[
t=\cos x.
\]

Получаем:
\[
\cos x=1.
\]

Общее решение:
\[
x=2\pi n,
\qquad n\in\mathbb Z.
\]

\textbf{Почему период \(2\pi\):} значение \(\cos x=1\) достигается после каждого полного оборота.

% -------------------------------------------------
\section*{\color{darkaccent}7. Отбор корней на промежутке}
\sectionline

Пусть требуется выбрать корни
\[
x=2\pi n
\]
на промежутке
\[
-\frac{7\pi}{2}\le x\le-2\pi.
\]

Подставляем общий корень в двойное неравенство:
\[
-\frac{7\pi}{2}
\le
2\pi n
\le
-2\pi.
\]

Делим все части на \(2\pi>0\):
\[
-\frac74\le n\le-1.
\]

Поскольку
\[
n\in\mathbb Z,
\]
получаем
\[
n=-1.
\]

Следовательно,
\[
x=-2\pi.
\]

\subsection*{\color{accent}Алгоритм отбора}

\begin{enumerate}
    \item Запишите заданный промежуток как двойное неравенство.
    \item Вместо \(x\) подставьте общий вид корня.
    \item Выразите целочисленный параметр \(n\) или \(k\).
    \item Выберите только целые значения параметра.
    \item Подставьте их в формулу корня.
\end{enumerate}

% -------------------------------------------------
\section*{\color{darkaccent}8. Ещё один пример с теоремой Безу}
\sectionline

Рассмотрим:
\[
3t^3+t^2-4t-2=0.
\]

Свободный член равен \(-2\), поэтому среди целых кандидатов:
\[
\pm1,\quad\pm2.
\]

Проверяем \(t=1\):
\[
3+1-4-2=-2\ne0.
\]

Проверяем \(t=-1\):
\[
-3+1+4-2=0.
\]

Значит,
\[
t=-1
\]
--- корень, и многочлен делится на
\[
t+1.
\]

После деления:
\[
3t^3+t^2-4t-2
=
(t+1)(3t^2-2t-2).
\]

Поэтому:
\[
(t+1)(3t^2-2t-2)=0.
\]

Корни:
\[
t=-1,
\]
а также
\[
3t^2-2t-2=0.
\]

\[
D=(-2)^2-4\cdot3\cdot(-2)=28,
\]

\[
t=\frac{2\pm\sqrt{28}}{6}
=
\frac{1\pm\sqrt7}{3}.
\]

Если эта замена возникла из
\[
t=\cos x
\quad\text{или}\quad
t=\sin x,
\]
после решения необходимо дополнительно проверить условие
\[
-1\le t\le1.
\]

% -------------------------------------------------
\section*{\color{darkaccent}9. Логарифмические неравенства: обязательный алгоритм}
\sectionline

Для выражения
\[
\log_a f(x)
\]
должны выполняться условия:
\[
a>0,\qquad a\ne1,\qquad f(x)>0.
\]

Если основание постоянно и корректно, основное внимание уделяется телам логарифмов.

Для сравнения
\[
\log_a f(x)\ge\log_a g(x)
\]
при
\[
a>1
\]
получаем:
\[
f(x)\ge g(x),
\]
при обязательных условиях
\[
f(x)>0,\qquad g(x)>0.
\]

Если
\[
0<a<1,
\]
знак неравенства меняется:
\[
\log_a f(x)\ge\log_a g(x)
\quad\Longleftrightarrow\quad
f(x)\le g(x).
\]

% -------------------------------------------------
\section*{\color{darkaccent}10. ОДЗ логарифмического неравенства}
\sectionline

Рассмотрим:
\[
\log_{11}(2x^2+1)
+
\log_{11}\left(\frac1{32x}+1\right)
\ge
\log_{11}\left(\frac{x}{16}+1\right).
\]

Основание:
\[
11>1,
\]
поэтому оно допустимо и функция логарифма возрастает.

ОДЗ определяется по \textbf{исходному} выражению.

Первое условие:
\[
2x^2+1>0.
\]

Оно выполняется при всех
\[
x\in\mathbb R.
\]

Второе:
\[
\frac1{32x}+1>0.
\]

Приводим к общему знаменателю:
\[
\frac{1+32x}{32x}>0.
\]

Критические точки:
\[
x=-\frac1{32},
\qquad
x=0.
\]

Отсюда:
\[
x\in
\left(-\infty;-\frac1{32}\right)
\cup
(0;+\infty).
\]

Третье:
\[
\frac{x}{16}+1>0,
\]
то есть
\[
x>-16.
\]

Пересечение всех условий:
\[
\boxed{
x\in
\left(-16;-\frac1{32}\right)
\cup
(0;+\infty)
}.
\]

% -------------------------------------------------
\section*{\color{darkaccent}11. Метод интервалов для дробного неравенства}
\sectionline

Для неравенства вида
\[
\frac{f(x)}{g(x)}>0
\]
нельзя просто умножать обе части на выражение \(g(x)\), если его знак заранее неизвестен.

Причины:

\begin{itemize}
    \item \(g(x)\) может обращаться в ноль;
    \item \(g(x)\) может быть положительным или отрицательным;
    \item при умножении на отрицательное выражение знак неравенства должен измениться.
\end{itemize}

Надёжный путь:

\begin{enumerate}
    \item перенести всё в одну сторону;
    \item привести к одной дроби;
    \item найти нули числителя;
    \item найти нули знаменателя;
    \item отметить их на числовой прямой;
    \item определить знаки на промежутках;
    \item выбрать нужные интервалы.
\end{enumerate}

% -------------------------------------------------
\section*{\color{darkaccent}12. Решение основного логарифмического неравенства}
\sectionline

Используем свойство:
\[
\log_a u+\log_a v=\log_a(uv),
\]
которое допустимо на ОДЗ.

Получаем:
\[
\log_{11}
\left[
(2x^2+1)
\left(
\frac1{32x}+1
\right)
\right]
\ge
\log_{11}
\left(
\frac{x}{16}+1
\right).
\]

Так как
\[
11>1,
\]
сравниваем тела:
\[
(2x^2+1)
\left(
\frac1{32x}+1
\right)
\ge
\frac{x}{16}+1.
\]

Приводим дроби к удобному виду:
\[
\frac{(2x^2+1)(1+32x)}{32x}
-
\frac{x+16}{16}
\ge0.
\]

Общий знаменатель:
\[
32x.
\]

Получаем:
\[
\frac{
(2x^2+1)(1+32x)-2x(x+16)
}{32x}
\ge0.
\]

Раскрываем скобки:
\[
(2x^2+1)(1+32x)
=
2x^2+64x^3+1+32x.
\]

Следовательно,
\[
2x^2+64x^3+1+32x
-
2x^2
-
32x
=
64x^3+1.
\]

Имеем:
\[
\frac{64x^3+1}{32x}\ge0.
\]

Числитель:
\[
64x^3+1=0
\quad\Longrightarrow\quad
x^3=-\frac1{64}
\quad\Longrightarrow\quad
x=-\frac14.
\]

Знаменатель:
\[
32x=0
\quad\Longrightarrow\quad
x=0.
\]

Точка \(x=0\) всегда исключается.

% -------------------------------------------------
\section*{\color{darkaccent}13. Числовая прямая}
\sectionline

Для
\[
\frac{64x^3+1}{32x}\ge0
\]
получаем критические точки
\[
-\frac14,\qquad 0.
\]

\begin{center}
\begin{tikzpicture}[x=4.2cm,y=1cm]
    \draw[-{Stealth[length=2mm]},line width=0.8pt]
        (-1.1,0) -- (1.1,0);

    \draw[line width=0.8pt]
        (-0.45,-0.08) -- (-0.45,0.08);
    \fill[accent] (-0.45,0) circle (2pt);
    \node[below=5pt] at (-0.45,0) {$-\frac14$};

    \draw[line width=0.8pt]
        (0.35,-0.08) -- (0.35,0.08);
    \draw[fill=white,line width=0.8pt]
        (0.35,0) circle (2pt);
    \node[below=5pt] at (0.35,0) {$0$};

    \node[above=5pt] at (-0.78,0) {$+$};
    \node[above=5pt] at (-0.05,0) {$-$};
    \node[above=5pt] at (0.73,0) {$+$};
\end{tikzpicture}
\end{center}

Следовательно, без учёта ОДЗ:
\[
x\in
\left(-\infty;-\frac14\right]
\cup
(0;+\infty).
\]

Теперь пересекаем это множество с ОДЗ:
\[
\left(-16;-\frac1{32}\right)
\cup
(0;+\infty).
\]

Получаем окончательный ответ:
\[
\boxed{
x\in
\left(-16;-\frac14\right]
\cup
(0;+\infty)
}.
\]

% -------------------------------------------------
\section*{\color{darkaccent}14. Кубический корень и знак числа}
\sectionline

Для нечётной степени отрицательное число допустимо:

\[
\sqrt[3]{-a}=-\sqrt[3]{a}.
\]

Например:
\[
x^3=-\frac1{64}
\quad\Longrightarrow\quad
x=-\frac14.
\]

Для корня чётной степени в действительных числах:
\[
\sqrt{-a}
\]
при \(a>0\) смысла не имеет.

Важно различать:
\[
x^2=-4
\]
--- действительных решений нет;

\[
x^3=-8
\]
--- имеет решение
\[
x=-2.
\]

% -------------------------------------------------
\section*{\color{darkaccent}15. Типичные ошибки}
\sectionline

\begin{itemize}
    \item \textbf{Забыть ограничение после замены}
    \[
    t=\sin x
    \quad\text{или}\quad
    t=\cos x.
    \]

    \item \textbf{Проверять ОДЗ уже преобразованного выражения.}
    ОДЗ определяется по исходной записи.

    \item \textbf{Умножать неравенство на выражение с переменной}, знак которого неизвестен.

    \item \textbf{Включать корень знаменателя} в ответ.

    \item \textbf{Забывать менять знак} при умножении или делении неравенства на отрицательное число.

    \item \textbf{Терять скобки при вычитании многочленов.}

    \item \textbf{Путать делитель свободного члена с гарантированным корнем.}
    Делители --- только кандидаты.

    \item \textbf{После нахождения одного корня кубического многочлена останавливаться.}
    Нужно выделить линейный множитель и решить оставшееся квадратное уравнение.

    \item \textbf{Не проверять дискриминант} оставшегося квадратного множителя.

    \item \textbf{Отбирать тригонометрические корни ``на глаз''}, когда надёжнее записать двойное неравенство.
\end{itemize}

% -------------------------------------------------
\section*{\color{darkaccent}16. Краткий алгоритм перед экзаменом}
\sectionline

\begin{enumerate}
    \item \textbf{Тригонометрия:}
    унифицируйте аргументы и функции.

    \item \textbf{Замена:}
    при \(t=\sin x\) или \(t=\cos x\) сразу запишите
    \[
    t\in[-1;1].
    \]

    \item \textbf{Кубический многочлен:}
    проверьте простые целые делители свободного члена.

    \item \textbf{Найден корень \(a\):}
    используйте множитель
    \[
    t-a.
    \]

    \item \textbf{Разделите многочлен}
    и решите оставшийся квадратный множитель.

    \item \textbf{Вернитесь к тригонометрической переменной}
    и получите общий вид \(x\).

    \item \textbf{Для промежутка}
    проведите отбор через двойное неравенство.

    \item \textbf{Для логарифмов}
    сначала выпишите ОДЗ.

    \item \textbf{Дробное неравенство}
    приведите к одной дроби и примените метод интервалов.

    \item \textbf{Перед ответом}
    пересеките решение преобразованного неравенства с ОДЗ.
\end{enumerate}

% =================================================
% ТРЕНИРОВОЧНЫЙ БЛОК
% =================================================
\clearpage

\section*{\color{darkaccent}Тренировочный блок}
\sectionline

\textbf{Уровни:} 1--3 --- средние; 4--9 --- выше среднего; 10 --- сложная.

Решения в этом разделе не приводятся. Рекомендуется оформлять все существенные преобразования и ограничения.

\begin{enumerate}

    \item \textbf{Средняя.}
    Решите уравнение
    \[
    2\cos^2x+\sin^2x=1.
    \]

    \item \textbf{Средняя.}
    Разложите на множители многочлен
    \[
    t^3-2t^2-t+2,
    \]
    используя подбор целого корня и теорему Безу.

    \item \textbf{Средняя.}
    Решите неравенство
    \[
    \log_3(x+2)\ge\log_3(5-x).
    \]

    \item \textbf{Выше среднего.}
    Решите уравнение
    \[
    2\cos^3x-\cos^2x-1=0.
    \]

    \item \textbf{Выше среднего.}
    Найдите все корни уравнения
    \[
    3t^3+t^2-4t-2=0.
    \]

    \item \textbf{Выше среднего.}
    Решите уравнение
    \[
    4\cos^3x-\cos^2x-\cos x=0.
    \]

    \item \textbf{Выше среднего.}
    Из общего решения
    \[
    x=\frac{\pi}{6}+\pi k,
    \qquad k\in\mathbb Z,
    \]
    выберите все корни, принадлежащие промежутку
    \[
    -\frac{5\pi}{2}\le x\le\frac{7\pi}{3}.
    \]

    \item \textbf{Выше среднего.}
    Решите неравенство
    \[
    \log_5(x+4)+\log_5(x-1)
    \ge
    \log_5(2x+6).
    \]

    \item \textbf{Выше среднего.}
    Решите неравенство
    \[
    \log_{1/2}(x+3)
    \le
    \log_{1/2}(7-x).
    \]

    \item \textbf{Сложная.}
    Решите неравенство
    \[
    \log_{11}(2x^2+1)
    +
    \log_{11}\left(\frac1{8x}+1\right)
    \ge
    \log_{11}\left(\frac{x}{4}+1\right).
    \]

\end{enumerate}

% =================================================
% ОТВЕТЫ
% =================================================
\clearpage

\section*{\color{darkaccent}Ответы к тренировочному блоку}
\sectionline

\renewcommand{\arraystretch}{1.35}

\begin{table}[ht]
\centering
\caption{Краткие ответы}
\begin{tabularx}{\linewidth}{@{}cX@{}}
\toprule
\textbf{№} & \textbf{Ответ} \\
\midrule

1 &
\[
x=\frac{\pi}{2}+\pi n,\qquad n\in\mathbb Z.
\]
\\

2 &
\[
(t-1)(t-2)(t+1).
\]
\\

3 &
\[
x\in\left[\frac32;5\right).
\]
\\

4 &
\[
x=2\pi n,\qquad n\in\mathbb Z.
\]
\\

5 &
\[
t=-1,\qquad
t=\frac{1-\sqrt7}{3},\qquad
t=\frac{1+\sqrt7}{3}.
\]
\\

6 &
\[
x=\frac{\pi}{2}+\pi n,
\]
или
\[
x=\pm\arccos\frac{1+\sqrt{17}}8+2\pi n,
\qquad n\in\mathbb Z.
\]
\\

7 &
\[
x\in
\left\{
-\frac{11\pi}{6},
-\frac{5\pi}{6},
\frac{\pi}{6},
\frac{7\pi}{6}
\right\}.
\]
\\

8 &
\[
x\in
\left[
\frac{1+\sqrt{41}}2;
+\infty
\right).
\]
\\

9 &
\[
x\in[2;7).
\]
\\

10 &
\[
x\in
\left(-4;-\frac12\right]
\cup
(0;+\infty).
\]
\\

\bottomrule
\end{tabularx}
\end{table}

\vfill

\begin{center}
\color{textgray}\small
Перед сверкой ответа проверьте ОДЗ, знаки на интервалах
и допустимость значений после тригонометрической замены.
\end{center}

\end{document}